\section{Introduction}

Direct-to-satellite IoT (D2S IoT) is emerging as a key connectivity paradigm for
low-power wide-area devices in areas lacking terrestrial infrastructure
\cite{todo_ntn_iot}.  Unlike conventional cellular or GNSS-assisted positioning,
a D2S terminal must operate with minimal prior knowledge of its own state: its
position, timing reference, and oscillator quality are all uncertain.  This
uncertainty complicates uplink scheduling, link adaptation, and any downstream
location-aware services.

A LEO satellite's high relative velocity with respect to a ground terminal
produces strong and structured Doppler dynamics.  The Doppler shift varies
predictably with geometry: as the satellite approaches, the observed carrier
frequency rises above the nominal value; as it recedes, it falls below.  This
temporal Doppler curve---the \emph{Doppler-time fingerprint}---is not merely an
impairment; it carries geometric information about the terminal's location
relative to the satellite ground track.

Estimating position from this fingerprint is challenging because the terminal's
low-quality oscillator introduces an unknown carrier frequency offset (CFO) and
drift, the terminal timing reference is offset from true UTC, and observations
are corrupted by noise.  These nuisance parameters are tightly coupled with
position in the observed frequency data, and single-time observations provide
insufficient geometric strength.

This paper makes the following contributions:

\begin{itemize}
  \item A Doppler-time fingerprint formulation that treats the full pass
        observation sequence as a geometric signature, explicitly separating
        it from oscillator nuisances.
  \item A nuisance-aware Bayesian grid posterior estimator that jointly
        infers coarse position, time offset, CFO, and drift over a bounded
        region of interest (ROI), with explicit marginalization over nuisance
        dimensions.
  \item A reproducible synthetic-pass and Monte Carlo evaluation pipeline
        that produces posterior heatmaps, error statistics, CRLB sensitivity
        diagnostics, and a Doppler-time ambiguity ablation study demonstrating
        why multi-epoch fingerprinting outperforms single-snapshot Doppler,
        all without requiring satellite over-the-air data.
  \item A conservative trace-driven hardware-in-the-loop (HIL) validation plan
        using an LR1121/STM32 packet source and a USRP B210 capture front-end,
        enabling validation of the signal-processing and observation-extraction
        pipeline in a lab setting without OTA claims.
\end{itemize}

No reinforcement-learning, PGRL, or semantic-communication framework is used.
All results are model-driven and reproducible from the provided scripts.
\end{section}